% Generated draft: requires mathematical and source review.
\documentclass[11pt]{article}
\usepackage[margin=25mm]{geometry}
\usepackage{fontspec}
\setmainfont{Latin Modern Roman}
\usepackage{amsmath,amssymb,mathtools,unicode-math}
\setmathfont{Latin Modern Math}
\usepackage{enumitem,needspace,xurl,hyperref}
\setlength{\parindent}{0pt}
\setlength{\parskip}{5pt}
\setlength{\emergencystretch}{3em}
\newcommand{\N}{\mathbb{N}}
\newcommand{\Z}{\mathbb{Z}}
\newcommand{\Q}{\mathbb{Q}}
\newcommand{\R}{\mathbb{R}}
\newcommand{\C}{\mathbb{C}}
\newcommand{\F}{\mathbb{F}}
\newcommand{\A}{\mathbb{A}}
\newcommand{\PP}{\mathbb{P}}
\newcommand{\E}{\mathbb{E}}
\newcommand{\eps}{\varepsilon}
\newcommand{\dd}{\,\mathrm{d}}
\begin{document}
\section*{UnsolvedMath ID: 10400146}
\textbf{Source problem code:} catalogue entry
\subsection{Definitions and mathematical statement}
An integral homology \(3\)-sphere is a closed oriented \(3\)-manifold \(M\) with \(H_i(M;\Z)\cong H_i(S^3;\Z)\) for every \(i\). Problem 7.31 asks to characterize which elements in the source's Habiro expansion (45) occur as \(\tau^{\mathrm{SO}(3)}(M)\) for such \(M\).

For terminology, the usual Habiro completion is
\[
 \widehat{\Z[q]}=\varprojlim_{n\ge1}
 \Z[q]\big/\left(\prod_{j=1}^{n}(1-q^j)\right).
\]
Its elements can be represented by convergent cyclotomic series of the form
\[
 \sum_{n\ge0} f_n(q)\prod_{j=1}^{n}(1-q^j),\qquad f_n(q)\in\Z[q],
\]
with the empty product equal to \(1\); such expressions need not be unique. This is a conventional explanation of Habiro expansion, not an identification of the missing equation (45) with a particular normalization or coefficient ring.

The requested characterization concerns realizability by actual integral homology \(3\)-spheres, rather than arbitrary elements or the linear span of realizable invariants. To state its exact target one must recover equation (45), the definition and normalization of \(\tau^{\mathrm{SO}(3)}\), any root-of-unity restrictions, and whether ``elements'' means elements of a completion or particular normalized expansion coefficients. The record supplies none of these details and gives no proposed necessary-and-sufficient condition.
\subsection{Short English statement}
Determine exactly which source-defined Habiro expansions of the \(\mathrm{SO}(3)\) quantum invariant arise from integral homology three-spheres. The referenced expansion (45) is missing, so its target and normalization remain unresolved. The supplied open-triage assessment reports no clear resolution.
\subsection{Sources}
\begin{itemize}
\item Local UnsolvedMath registry, permanent ID \(10400146\), exact source code \texttt{catalogue entry}, Problem 7.31. The statement refers to Habiro's expansion (45), not included in the record. No individual proposer or proposed year is supplied.
\item Ohtsuki, \emph{Problems on invariants of knots and 3-manifolds} (2002), Problem 7.31, PDF page 120: \url{https://msp.org/gtm/2002/04/gtm-2002-04-024s.pdf}. Bold printed headings distinguish 231 numbered items from surrounding references.
\item Extraction pdf-font-structure, accessed 2026-07-29. Source evidence NEEDS\_REVIEW presents an open problem and requires a current primary-literature audit. Supplied classification OPEN-TRIAGE calls characterization of the Habiro expansion not obviously resolved; no separate primary citation for (45) is supplied.
\item Dataset attribution: Ulam AI / UnsolvedMath Contributors, CC BY 4.0. Source rights NEEDS\_REVIEW: publicly accessible author/publisher copy, redistribution terms requiring review.
\end{itemize}
\end{document}
